% ---------------------------------------------------------------
\section{Spectral Ranking}
\label{sec:spectral_ranking}
% ---------------------------------------------------------------

\subsection{Units and corpus}

The ranking procedure operates over two disjoint sets of units: a set
$\mathcal{S}$ of \emph{sources} (academic journals and other publication
venues) and a set $\mathcal{U}$ of \emph{institutions}.  Both sets are ranked
simultaneously from the same citation data; the field of study determines
which works enter the corpus but does not alter the ranking mechanism.

For a given field $f$ and census window $[t_0, t_1]$, a work is included in
the corpus if it (i) has at least one topic tag associated with field $f$ in
the OpenAlex taxonomy, (ii) was published during $[t_0, t_1]$, and (iii) is
affiliated with at least one retained institution (defined below).  Throughout
this study the window is 2020--2024 and the field taxonomy follows the 26
top-level OpenAlex fields.

\subsection{Candidacy and retention thresholds}

Not every source or institution that appears in the corpus is ranked.  Units
must clear a minimum activity threshold $\tau$ expressed in
\emph{weighted works per year}: a source $s$ is retained if

\[
  \tilde{w}_s \;=\; \frac{1}{T}\sum_{t=t_0}^{t_1} \sum_{\substack{p \in \mathcal{P}_f \\ \mathrm{source}(p)=s}} f_{p,f}
  \;\geq\; \tau_s,
\]

where $T = t_1 - t_0 + 1$ is the window length in years, $\mathcal{P}_f$ is
the set of corpus works for field $f$, and $f_{p,f} \in (0,1]$ is the
fractional topic score of work $p$ in field $f$ (a work may straddle
multiple fields, with scores summing to at most one per work).  An analogous
threshold $\tau_u$ applies to institutions, with each work's fractional field
weight allocated to each affiliated institution in proportion to that
institution's author share.  The baseline threshold is $\tau_s = \tau_u =
10$~weighted works per year, yielding a minimum of 50 weighted works over the
five-year window.  Units that fall below their threshold are excluded from the
ranking but still participate as citers and cited works in the edge list; they
do not receive a prestige score.

\subsection{Citation edge list}

The edge list is derived from the full set of within-window citation pairs
$(p, q)$ where both $p$ (the citing work) and $q$ (the cited work) belong to
the retained corpus for field $f$.  Each such pair generates one or more
directed edges in the bipartite graph according to the Cartesian product of
the institutions affiliated with $p$ and those affiliated with $q$.

Each edge $(p, q, u_p, u_q)$ — citer work $p$, cited work $q$, citer
institution $u_p$, cited institution $u_q$ — carries a weight composed of
three multiplicative factors.  The \emph{field weight} $f_{p,f}$ scales each
citer's contribution by its fractional relevance to the field.  The
\emph{institution weight} $w_{u_p}$ distributes the citer's credit across its
affiliated institutions according to the author--institution share (the
fraction of authors of $p$ affiliated with $u_p$); analogously the
\emph{cited institution weight} $w_{u_q}$ distributes cited credit across
institutions of $q$.  An optional \emph{reference-count normalisation} factor
$\rho_p = \bar{R}/R_p$ (where $R_p$ is the number of references in work $p$
and $\bar{R}$ is the mean across all citing works) ensures that high-volume
citers do not dominate the flow matrix.  In the baseline configuration ($\rho
= 0$) this normalisation is applied; setting $\rho = 1$ uses raw reference
counts without normalisation.

\subsection{Bipartite citation-flow matrices}

Two raw citation-flow matrices are assembled from the edge list.  Let $N_s =
|\mathcal{S}|$ and $N_u = |\mathcal{U}|$ denote the numbers of retained
sources and institutions respectively.  The \emph{source-to-institution} matrix
$C_{SI} \in \mathbb{R}^{N_s \times N_u}$ accumulates the weighted citation
mass flowing from each source to each institution:
\begin{align}
  C_{SI}[s, u'] &= \sum_{\substack{(p,q,u_p,u_q) \in \mathcal{E} \\
                   \mathrm{src}(p)=s,\;u_q=u'}} \!\!
                   f_{p,f} \cdot \rho_p \cdot w_{u'},
  \label{eq:csi}
\end{align}
where the sum runs over all edge records in $\mathcal{E}$ whose citer work $p$
is published in source $s$ and whose cited work $q$ is affiliated with
institution $u'$.  The complementary \emph{institution-to-source} matrix
$C_{IS} \in \mathbb{R}^{N_u \times N_s}$ accumulates the reverse flow:
\begin{align}
  C_{IS}[u, s'] &= \sum_{\substack{(p,q,u_p,u_q) \in \mathcal{E} \\
                   u_p=u,\;\mathrm{src}(q)=s'}} \!\!
                   f_{p,f} \cdot \rho_p \cdot w_u.
  \label{eq:cis}
\end{align}

Together $C_{SI}$ and $C_{IS}$ encode all citation information needed for the
bipartite ranking: prestige propagates from sources to institutions via $C_{SI}$
and back from institutions to sources via $C_{IS}$, so that sources are ranked
by the quality of the institutions that cite them and institutions are ranked by
the quality of the sources they cite.

\subsection{Row normalisation and transition matrices}

Each matrix is row-normalised to obtain a stochastic transition matrix.  For
$C_{SI}$, row $s$ has outgoing weight $r_s = \sum_{u'} C_{SI}[s, u']$; where
$r_s > 0$ we set $H_{SI}[s,\cdot] = C_{SI}[s,\cdot] / r_s$, and analogously
for $H_{IS}$.  Rows with $r_x = 0$ (units that cite nothing within the retained
corpus) are left as zero rows and treated as \emph{dangling} nodes: at each
iteration of the power method their probability mass is redistributed through
the prior, exactly as in the standard dangling-node treatment for PageRank.

\subsection{Bipartite spectral ranking}

The two stochastic matrices $H_{SI} \in \mathbb{R}^{N_s \times N_u}$ and
$H_{IS} \in \mathbb{R}^{N_u \times N_s}$ are coupled: a
two-step random walk that alternates source$\to$institution$\to$source defines
the \emph{one-mode projection} onto sources:
\[
  M_S \;=\; H_{SI}\, H_{IS} \;\in\; \mathbb{R}^{N_s \times N_s}.
\]

The source prestige vector $\boldsymbol{\pi}_S$ is the Perron (left)
eigenvector of $M_S$:
\[
  M_S^\top \boldsymbol{\pi}_S \;=\; \lambda_1 \boldsymbol{\pi}_S,
  \qquad \lambda_1 = 1, \quad \boldsymbol{\pi}_S \geq 0, \quad
  \|\boldsymbol{\pi}_S\|_1 = 1.
\]

Because $M_S$ is a product of row-stochastic matrices it is itself
substochastic (or stochastic when all rows of $H_{SI}$ and $H_{IS}$ are
non-zero); under mild primitivity conditions guaranteed by the retained
citation graph, the Perron--Frobenius theorem asserts the existence and
uniqueness of $\boldsymbol{\pi}_S$.  In practice the eigenvector is computed
via ARPACK's implicitly restarted Lanczos method applied to $M_S^\top$,
requesting the two dominant eigenpairs so that both $\lambda_1$ and $\lambda_2$
are recovered simultaneously.

Institution prestige scores are recovered by one reverse hop:
\[
  \boldsymbol{\pi}_U \;=\; H_{SI}^\top \boldsymbol{\pi}_S,
\]
and then independently $\ell_1$-normalised.  This ensures that
$\|\boldsymbol{\pi}_S\|_1 = \|\boldsymbol{\pi}_U\|_1 = 1$ and that the two
scales are comparable.

\subsection{Prestige scores}

The raw eigenvector components $\pi_s$ and $\pi_u$ are proportional to the
stationary probability of the random walk arriving at each unit and are
therefore sensitive to corpus size: a unit in a large field accumulates more
probability mass than an equally prestigious unit in a small field, simply
because the walk has more mass to distribute.  To obtain a scale-invariant
prestige score we normalise by each unit's share of the corpus activity $a_p$
(its weighted works over the window):
\[
  v_p \;=\; \frac{A \cdot \pi_p}{A_{\mathrm{real}} \cdot a_p},
\]
where $A = \sum_{p \in \mathcal{S} \cup \mathcal{U}} a_p$ is the total
corpus activity and $A_{\mathrm{real}}$ is the sum of $\pi_p$ over
non-sentinel units (relevant only when cross-field sentinel edges are used).
Under this normalisation $v_p$ has the interpretation of prestige per unit of
activity: a unit with $v_p = 1$ receives exactly the corpus-average citation
return for its level of output, units with $v_p > 1$ are above average, and
units with $v_p < 1$ are below average.  The prestige score is the primary
output used throughout this paper.

\subsection{Work-level prestige}

To support the enclave analysis, we require a prestige score at the level of
individual works rather than venues.  For a work $p$ published in source $s$
and affiliated with institutions $\{u_1, \ldots, u_k\} \cap \mathcal{U}$ (the
subset that are retained), we define
\[
  v_p^{\mathrm{work}} \;=\;
  \frac{v_s + \overline{v}_u}{2},
\]
where $v_s$ is the prestige score of the source and $\overline{v}_u =
\frac{1}{k}\sum_{i=1}^{k} v_{u_i}$ is the unweighted mean over retained
affiliated institutions.  When no affiliated institution is retained
($k = 0$), the source score alone is used: $v_p^{\mathrm{work}} = v_s$.  The
equal 50--50 split between source and institutional prestige reflects a design
choice that neither venue reputation nor institutional reputation dominates; it
is consistent across all 26 fields and the full five-year window.

\subsection{Spectral gap and field structure}

A by-product of requesting two ARPACK eigenpairs is the spectral gap
$\Delta = 1 - |\lambda_2|$, where $\lambda_2$ is the subdominant eigenvalue
of $M_S$.  A large gap (close to 1) indicates that the Perron eigenvector
captures a strongly unified hierarchy: prestige flows coherently across the
field and the random walk mixes rapidly to its stationary distribution.  A
small gap indicates near-reducibility: the field decomposes into weakly
connected sub-communities that exchange little citation mass, so that the
second eigenvector carries substantial structural information about the
partition of the field.

Across the 26 OpenAlex fields the gap ranges from 0.091 (Business, Management
and Accounting) to 0.860 (Neuroscience).  Social science and humanities fields
tend toward small gaps, reflecting the plurality of paradigms and citation
communities within those broad umbrellas.  Physical and biomedical science
fields tend toward large gaps, consistent with more unified citation hierarchies.
The gap is a diagnostic, not a quality indicator; the enclave analysis applies
identically regardless of gap size.

\subsection{Parameter summary}

Table~\ref{tab:params} summarises the parameters of the ranking procedure and
their baseline values used throughout this paper.

\begin{table}[ht]
\centering
\caption{Ranking parameters and baseline values.}
\label{tab:params}
\begin{tabular}{lllp{6.2cm}}
\toprule
Symbol & Baseline & Range & Description \\
\midrule
$\tau_s$, $\tau_u$ & 10/yr & $>0$ & Retention threshold (weighted works per year) \\
$\rho$ & 0 & $\{0,1\}$ & Reference-count normalisation; 0 = fixed-count ($\bar{R}/R_i$), 1 = full-count \\
$\omega$ & 0 & $\{0,1\}$ & Institution weighting; 0 = author-fractional, 1 = uniform $1/N_{\mathrm{inst}}$ \\
$\alpha$ & 1 & $(0,1]$ & Katz damping; $1$ = pure Perron eigenvector \\
\bottomrule
\end{tabular}
\end{table}
